%!TEX root =  tfg.tex
\chapter{Iteración 1: Autenticación OAuth2 con Google}

\begin{abstract}
En esta primera iteración se implementa el sistema de autenticación de usuarios mediante OAuth2 con Google. El objetivo es permitir que los usuarios inicien sesión con su cuenta de Google, obtener y almacenar su información básica de perfil (nombre, email y foto), y proteger los recursos de la API con un token JWT propio emitido por el backend.
\end{abstract}

\section{Características a desarrollar}

\begin{enumerate}
\item Autenticación OAuth2 con Google (flujo de ID Token). Ver Tabla \ref{tab:valorAportado1}.
\item Emisión y validación de JWT propio desde el backend.
\item Página de perfil de usuario con nombre y foto.
\item Almacenamiento del token JWT en el frontend (localStorage + MobX).
\end{enumerate}

\begin{table*}[htb]
	\centering
	\begin{coolTable}{p{4cm}p{\textwidth-4.5cm}}{2}
{Análisis de valor aportado 0001}
\textbf{Propuesta}&Autenticación de usuarios mediante Google OAuth2\\
\midrule
\textbf{Valor}&Permite identificar a los usuarios de forma segura sin gestionar contraseñas. Usa la infraestructura de Google, lo que simplifica el flujo de registro y evita almacenar credenciales sensibles.\\
\textbf{Coste}&Requiere configurar un proyecto en Google Cloud Console, añadir la librería \texttt{@react-oauth/google} al frontend y JJWT al backend. Coste en esfuerzo: bajo. Coste económico: nulo (servicio gratuito).\\
\textbf{Opciones}&Otras opciones consideradas: Auth0 (servicio de terceros, introduce dependencia externa), Keycloak (alta complejidad de despliegue), gestión propia de usuarios con contraseñas (mayor superficie de ataque).\\
\textbf{Riesgos}&Cambios en la API de Google. Expiración del token de Google en frontend. La URL de callback debe actualizarse en Google Cloud al cambiar de dominio.\\
\textbf{Deuda técnica}&El token JWT se almacena en localStorage, lo que es susceptible a ataques XSS. Una mejora futura sería usar cookies HttpOnly.\\
	\end{coolTable}
	\caption{Análisis de valor aportado 0001\label{tab:valorAportado1}}
\end{table*}

\section{Diseño}

El flujo de autenticación implementado es el siguiente:

\begin{enumerate}
    \item El usuario pulsa el botón ``Iniciar sesión con Google'' en la página principal.
    \item El frontend obtiene un \textit{ID Token} de Google mediante la librería \texttt{@react-oauth/google}.
    \item El frontend envía el ID Token al endpoint \texttt{POST /api/auth/google} del backend.
    \item El backend valida el token llamando al endpoint público de Google \texttt{tokeninfo} y verifica que el \textit{audience} coincide con el Client ID de la aplicación.
    \item Si el token es válido, el backend crea o recupera el usuario en base de datos y genera un JWT propio firmado con HMAC-SHA256.
    \item El frontend almacena el JWT y los datos del usuario en el \texttt{UserStore} (MobX) y en \texttt{localStorage} para persistencia entre sesiones.
\end{enumerate}

\begin{figure}[htbp]
\begin{center}
\missingfigure{Diagrama de secuencia: flujo OAuth2 Google + emisión JWT propio}
\caption{Diagrama de flujo de autenticación OAuth2 con Google}
\label{fig:diseno01}
\end{center}
\end{figure}

\begin{table*}[htb]
	\centering
	\begin{coolTable}{p{4cm}p{\textwidth-4.5cm}}{2}
{Memorando técnico 0001}
\textbf{Asunto}&Elección del flujo de autenticación: Frontend-driven vs. Backend-driven\\
\textbf{Resumen}&Se utiliza el flujo \textit{frontend-driven}: el frontend obtiene el ID Token de Google y lo envía al backend para validación, en lugar del flujo tradicional con redirecciones OAuth2 server-side.\\
\midrule
\textbf{Factores causantes}&La aplicación es una SPA desplegada en Vercel con un backend desacoplado en Railway. El flujo server-side con redirecciones requeriría configurar correctamente las URIs de callback en ambas plataformas y gestionar cookies cross-domain, lo que introduce complejidad adicional.\\
\textbf{Solución}&El frontend usa \texttt{GoogleLogin} de \texttt{@react-oauth/google}, que devuelve un ID Token (JWT firmado por Google). Este token se envía al backend, que lo valida contra el endpoint \texttt{https://oauth2.googleapis.com/tokeninfo} y emite un JWT propio.\\
\textbf{Motivación}&Simplicidad de implementación, compatibilidad con arquitecturas SPA + API REST, y eliminación de problemas de CORS con cookies en dominios distintos.\\
\textbf{Cuestiones abiertas}&El almacenamiento del JWT en \texttt{localStorage} es vulnerable a XSS. Se podría mejorar con cookies HttpOnly en una iteración futura.\\
\textbf{Alternativas}&Spring Security OAuth2 Resource Server (más complejo, requiere integración Spring Security completa); Auth0 (introduce dependencia de terceros y coste potencial).\\
	\end{coolTable}
	\caption{Memorando técnico 0001}
\end{table*}

\section{Implementación}

\subsection{Backend}

Se añaden las siguientes clases al proyecto Spring Boot:

\begin{itemize}
    \item \texttt{User}: entidad JPA que almacena \texttt{googleId}, \texttt{email}, \texttt{name} y \texttt{pictureUrl}.
    \item \texttt{UserRepository}: repositorio Spring Data JPA con método \texttt{findByGoogleId}.
    \item \texttt{UserService}: servicio que implementa el patrón \textit{find-or-create} al autenticar.
    \item \texttt{JwtService}: genera y valida tokens JWT firmados con HMAC-SHA256 usando la librería JJWT 0.12.6. El secreto se inyecta desde la variable de entorno \texttt{JWT\_SECRET}.
    \item \texttt{JwtAuthFilter}: filtro \texttt{OncePerRequestFilter} que extrae el token del header \texttt{Authorization: Bearer} y establece el contexto de seguridad de Spring.
    \item \texttt{AuthController}: expone \texttt{POST /api/auth/google}. Valida el ID Token con Google, crea el usuario y devuelve el JWT propio.
\end{itemize}

La configuración de seguridad (\texttt{SecurityConfig}) se actualiza para registrar el \texttt{JwtAuthFilter} antes de \texttt{UsernamePasswordAuthenticationFilter} y establecer la política de sesión como \texttt{STATELESS}.

\begin{asigResponsabilidad}{AUTH001}{AuthController}
{ResponseEntity<?> loginWithGoogle(Map<String,String> body)}
\pasoPseudo{1. Extraer campo \texttt{credential} del body.}
\pasoPseudo{2. Llamar a \texttt{https://oauth2.googleapis.com/tokeninfo?id\_token=\{credential\}}.}
\pasoPseudo{3. Verificar que el campo \texttt{aud} o \texttt{azp} coincide con el Client ID configurado.}
\pasoPseudo{4. Llamar a \texttt{UserService.findOrCreate()} con \texttt{sub}, \texttt{email}, \texttt{name}, \texttt{picture}.}
\pasoPseudo{5. Generar JWT con \texttt{JwtService.generateToken(user)}.}
\pasoPseudo{6. Devolver \texttt{\{token, user\}} con HTTP 200.}
\cabeceraMetodosBajoNivel
\pasoCodigo{2}{RestTemplate}{getForObject(url, Map.class)}{001}{SI}
\pasoCodigo{4}{UserService}{findOrCreate(googleId, email, name, picture)}{001}{SI}
\pasoCodigo{5}{JwtService}{generateToken(user)}{001}{SI}
\end{asigResponsabilidad}

\subsection{Frontend}

Se añaden los siguientes elementos:

\begin{itemize}
    \item \texttt{UserStore}: store MobX con campos \texttt{user} y métodos \texttt{setUser}, \texttt{logout}. Persiste en \texttt{localStorage} para mantener la sesión entre recargas.
    \item \texttt{RootStore}: se añade \texttt{userStore} junto a \texttt{speciesStore} y \texttt{mapStore}.
    \item \texttt{GoogleOAuthProvider}: envuelve la aplicación en \texttt{main.tsx} con el Client ID configurado en la variable de entorno \texttt{VITE\_GOOGLE\_CLIENT\_ID}.
    \item \texttt{Home.tsx}: se añade el botón de Google Login en la esquina superior derecha. Si el usuario ya está autenticado, se muestra su avatar y nombre con enlace al perfil.
    \item \texttt{Profile.tsx}: página de perfil que muestra foto, nombre y email del usuario, con opción de cerrar sesión.
\end{itemize}

\section{Pruebas}

Las pruebas de esta iteración se realizan de forma manual:

\begin{enumerate}
    \item Verificar que el botón de Google Login aparece en la Home cuando no hay sesión activa.
    \item Iniciar sesión con una cuenta Google válida y comprobar que aparece el avatar del usuario.
    \item Navegar a la página de perfil y verificar que se muestran nombre, email y foto.
    \item Cerrar sesión y comprobar que se borra el estado y se redirige a la Home.
    \item Recargar la página con sesión activa y verificar que se restaura desde \texttt{localStorage}.
    \item Enviar un token inválido al endpoint \texttt{/api/auth/google} y verificar que devuelve HTTP 401.
\end{enumerate}

\section{Despliegue}

Se añaden las siguientes variables de entorno en los entornos de producción y staging:

\begin{itemize}
    \item Backend (Railway): \texttt{GOOGLE\_CLIENT\_ID} y \texttt{JWT\_SECRET}.
    \item Frontend (Vercel): \texttt{VITE\_GOOGLE\_CLIENT\_ID}.
\end{itemize}

Los orígenes \texttt{https://scubex.vercel.app} y \texttt{https://scubex-git-develop-izshins-projects.vercel.app} se registran como orígenes autorizados en el proyecto de Google Cloud Console.


\begin{quotation}[Novelist]{Ernest Hemingway (1899--1961)}
The good parts of a book may be only something a writer is lucky enough to overhear or it may be the wreck of his whole damn life -- and one is as good as the other.
\end{quotation}

\begin{abstract}
Resumen de lo que va a ocurrir en el capítulo. ¿Cuál es el objetivo que tenemos con este capítulo?
\end{abstract}

\section{Características a desarrollar}

\begin{enumerate}
\item Funcionalidad A. Ver Tabla \ref{tab:valorAportado1}.
\item Funcionalidad B.
\end{enumerate}

\begin{table*}[htb]
	\centering
	\begin{coolTable}{p{4cm}p{\textwidth-4.5cm}}{2}
{Análisis de valor aportado 0001}
\textbf{Propuesta}&Trabajo que pretende analizarse y justificarse\\
\midrule
\textbf{Valor}&Qué valor aporta al proyecto o al usuario final.\\
\textbf{Coste}&Qué costes en términos de esfuerzo, adquisiciones y limitaciones tiene la propuesta\\
\textbf{Opciones}&Qué otras opciones se tienen que aporten un valor similar? ¿Es realmente un valor relevante para el proyecto/cliente\\
\textbf{Riesgos}&Qué riesgos pueden surgir a la hora de desarrollar esta propuesta.\\
\textbf{Deuda técnica}&Posibles deudas técnicas que se asumen con el desarrollo de esta propuesta.\\
	\end{coolTable}
	\caption{Análisis de valor aportado 0001\label{tab:valorAportado1}}
\end{table*}

\section{Diseño}
Aquí una discusión de cómo va a afectar todo al diseño

Debe insertarse un diagrama UML de diseño con los cambios y hacer referencia en el texto así Fig. \ref{fig:diseno01}.

\begin{figure}[htbp]
\begin{center}
\missingfigure{Aquí el modelo de diseño en formato vectorial preferentemente (pdf)}
% Incluir la figura quitando el comentario a la fila de abajo.
% \includegraphics[width=\textwidth]{myfile.pdf}
\caption{Diagrama UML de diseño para la iteración 1}
\label{fig:diseno01}
\end{center}
\end{figure}

Un memorando técnico por cada decisión de diseño.

\begin{table*}[htb]
	\centering
	\begin{coolTable}{p{4cm}p{\textwidth-4.5cm}}{2}
{Memorando técnico 0001}
\textbf{Asunto}&¿Cuál es el problema?\\
\textbf{Resumen}&¿Cuál es la solución propuesta?\\
\midrule
\textbf{Factores causantes}&Descripción pormenorizada del problema\\
\textbf{Solución}&Descripción pormenorizada de la solución propuesta\\
\textbf{Motivación}&¿Por qué propone esta solución?\\
\textbf{Cuestiones abiertas}& Factores a tener en cuenta en la solución cuya dimensión se reconoce.\\
\textbf{Alternativas}&Otras soluciones consideradas y la razón por la que se excluyeron.\\
	\end{coolTable}
	\caption{Memorando técnico 0001}
\end{table*}


\section{Implementación}

Un memorando técnico por cada decisión de implementación y refactorización que afecte al diseño.

\begin{asigResponsabilidad}{0001}{Prueba}
{[return\_type] method\_name1 (param1:type1, ...)}
\pasoPseudo{1. Paso 1.}
\pasoPseudo{2. Paso 2.}
\cabeceraMetodosBajoNivel
\pasoCodigo{1}{ClassName}{[return\_type] method\_name1 (param1:type1, ...)}{001}{SI}
\diagramaColaboracion{figures/colDiagram.png}
\end{asigResponsabilidad}

\begin{asigResponsabilidad}{alvotermar02}{Grubber}
{[return\_type] grubber (param1:type1, ...)}
\pasoPseudo{1. Lanzar 2 dados}
\pasoPseudo{2. Compara resultado de los dados con kicking del open-side}
\pasoPseudo{2.1. Si valor dados es menor o igual a kicking, avanza 10m}
\pasoPseudo{3.1. Si no hay defensa y el golpeo es exitoso, el pateador retiene la posesi\'on del bal\'on}
\pasoPseudo{3.2. Si hay defensa y el golpe es exitoso, el atacante tira un dado y suma su valor al de speed y strength y el defensor lanza 2 dados y lo suma al valor de speed y strength de su jugador, el vencedor ser\'a aquel que tenga m\'as puntos, si es igual, la posesi\'on es del defensor}
\pasoPseudo{4.1. Si no es exitoso y hay defensa el bal\'on pasa a posesi\'on del defensor}
\pasoPseudo{4.2. Si no es exitoso y no hay defensa de lanza un line-out}
\cabeceraMetodosBajoNivel
\pasoCodigo{1}{Dice}{[Integer] throwDice ()}{001}{SI}
\pasoCodigo{2}{ClassName}{[Int] compareKickingToDice (kicking:Integer, dice: Integer)}{001}{SI}
\pasoCodigo{2.1}{ClassName}{[Integer] setLine (line:Integer)}{001}{SI}
\pasoCodigo{4.2}{ClassName}{[Integer] lineOut ()}{001}{SI}
%\diagramaColaboracion{colDiagram.png}
\end{asigResponsabilidad}

\section{Pruebas}

Descripción de las pruebas realizadas al software

\section{Despliegue}

Breve resumen de cómo se han desplegado los cambios en el sistema de producción.